\begin{table}[H]
  \centering
  \caption{GPH estimates of the fractional differencing parameter $d$ for absolute returns of gold and silver}
  \label{tab:arfima-gph}
  \begin{tabular}{lcccccc}
    \toprule
    Asset & $m$ & $\hat{d}$ & s.e.($\hat{d}$) & $t$($d=0$) & $t$($d=0.5$) & $t$($d=1$) \\
    \midrule
    Gold & 68 & 0.576 & 0.090 & 6.39 (0.000) & 0.84 (0.399) & -4.70 (0.000) \\
    Silver & 68 & 0.552 & 0.080 & 6.91 (0.000) & 0.65 (0.515) & -5.60 (0.000) \\
    \bottomrule
  \end{tabular}
  \vspace{0.2cm}
  \begin{minipage}{0.9\textwidth}
  \footnotesize
  \textit{Notes}: The table reports Geweke--Porter--Hudak log-periodogram estimates of the fractional differencing parameter $d$ for the series of absolute daily log returns on GLD and SLV (in percentage points), using the first $m = \lfloor T^{0.5} \rfloor$ Fourier frequencies. Columns $t$($d=d_0$) report $(\hat{d}-d_0)/\mathrm{s.e.}(\hat{d})$ with two-sided normal $p$-values in parentheses. Values $0 < d < 0.5$ correspond to covariance-stationary long memory and $0.5 \le d < 1$ to non-stationary but mean-reverting long memory. Both point estimates exceed 0.5, i.e. they lie in the non-stationary, mean-reverting range, which is consistent with fractionally differencing the series before further modelling. At the 5\% level, $d=0$ is rejected for both assets; $d=0.5$ is not rejected for either asset; $d=1$ is rejected for both assets. This is a semiparametric fractional-differencing analysis; no parametric ARFIMA model is estimated.
  \end{minipage}
\end{table}
